\documentclass[11pt]{article}
\usepackage{fontspec}
\setmainfont{Times New Roman}
\setsansfont{Arial}
\setmonofont{Consolas}
\usepackage{amsmath,amssymb}
\usepackage{geometry}
\usepackage{microtype}
\geometry{a4paper,margin=3cm}
\setlength{\parindent}{1.25em}
\setlength{\parskip}{0pt}
\emergencystretch=10em
\allowdisplaybreaks
\providecommand{\Kfield}{\mathfrak{K}}
\providecommand{\Ssys}{\mathfrak{S}}
\providecommand{\Mfield}{\mathfrak{M}}
\providecommand{\tq}{\tau}
\providecommand{\ds}{\displaystyle}

\begin{document}

\begin{center}
{\Large\bfseries 5. Polja racionalnyh funkcij}\par
\vspace{1.2em}
J. Ber. d. DMV 22 (1913), S. 316--319
\end{center}

\vspace{2em}

Slědujuća pytanja iznačalno idut nazad k razgovoram s E. Fischerom. Něktore iz pytanj, vpročem, za specialny slučaj jednoj neopreděljenej -- i to pod obćejšimi prědpostavkami o koeficientnom polju -- už byly postavjene i rěšene E. Steinitzom.\footnote{E. Steinitz: Algebraische Theorie der Körper (osoblivo \S\ 24). Crelles Journal 137. 1910.}

Pod ``poljem racionalnyh funkcij'' razuměju polje, ktorego elementy sut racionalne funkcije od $n$ neopreděljenyh, s koeficientami iz danogo čislovogo polja, ktore osoblivo može obimati i vse kompleksne čisla. Priměry takyh polj sut polje simetričnyh funkcij od $n$ veličin, ili obćejše cělost racionalnyh funkcij od $n$ neopreděljenyh, ktore dopuščajut permutacije někojej grupy (Lagrangeve rodove oblasti); dalje, invariantno polje.

Medžu dvoma prvo nazvanymi poljami i poslědnim stoji harakterny rozděl: prva soderžat $n$ algebraično nezavisnyh funkcij, elementarne simetrične funkcije; dokolě čislo algebraično nezavisnyh invariantov jest vsegda menje než čislo neopreděljenyh. Zato jest važno, že obće možno taka polja drugogo tipa jednoznačno prirędžati poljam prvogo tipa, zaměnjajųči něktore iz neopreděljenyh čislami. Zato budu se v slědujućem ograničati na polja prvogo tipa, to jest na polja s $n$ algebraično nezavisnymi funkcijami; po toj prirędženosti rezultaty tada važet obće.

Pytanja grupujut se okolo treh bazovyh pojmov: racionalna baza, minimalna baza i integralna baza.

1. Pod ``racionalnoju bazoju'' razuměju konečno čislo funkcij polja tako, že vsaka funkcija polja može se predstaviti kako racionalna kombinacija togo konečnogo čisla funkcij, s koeficientami iz danogo koeficientnogo područja. Prostymi razmysljenjami -- ktore v slučaju jednoj neopreděljenej nahodet se takože u E. Steinitza, tamže -- dokazuje se eksistencija racionalnoj bazy za vsako polje racionalnyh funkcij. Napr., po Lagrangevoj teoremě elementarne simetrične funkcije i jedna funkcija, ktora naleži k grupě, tvoret racionalnu bazu za raněje spomenute Lagrangeve rodove oblasti.

2. Vsaka racionalna baza mora (za polja prvogo tipa) soderžati najmněje $n$ funkcij; ako ona soderži točno $n$ funkcij, ktore tada muset byti algebraično nezavisne, nazyvam ju ``minimalnoju bazoju''. Pytanje o eksistenciji minimalnoj bazy možno čestično rěšiti teoremami Lürotha, Castelnuova i Enriquesa.\footnote{J. Lüroth: Beweis eines Satzes über rationale Kurven. Math. Ann. 9. 1875. --- G. Castelnuovo: Sulla razionalità delle involuzioni piane. Math. Ann. 44. 1893. --- F. Enriques: Sopra una involuzione non razionale dello spazio. Rend. Acc. Linc. vol. 21. 21. Jan. 1912.} Po njih za polja od jednoj neopreděljenej vsegda eksistuje minimalna baza: Lürothova funkcija polja, jednoznačno opreděljena do drobno-linearnogo prěobraženja (sr. Steinitz \S\ 24). Za polja od dvoh neopreděljenyh minimalna baza eksistuje vsegda i obće samo togda, kogda dopušča se algebraično razširjenje koeficientnogo polja, dokolě za tri i bolje neopreděljenyh minimalna baza obće nikako ne eksistuje. Harakterizovati specialne polja s minimalnoju bazoju ješče ne bylo pokušanu.

Ja jesm nyně dalje slědila pytanje o minimalnoj bazě pri Lagrangevyh rodovyh oblastjah. Tu ono dobiva slědujuće značenje: ``Ako Lagrangeva rodova oblast, naležeča k grupě $G$, imaje minimalnu bazu, togda možno cělost jednačenj s prědpisanoju grupoju $G$ racionalno konstruovati črez parametrično predstavjenje; i to za vsako čislovo polje, ktore soderži koeficientno područje minimalnoj bazy -- znači osoblivo za vsako libovoljno čislovo polje, ako to koeficientno područje jest polje racionalnyh čisel.''

Najprostejši priměr davaje simetrična grupa; tu elementarne simetrične funkcije tvoret minimalnu bazu s racionalnymi čislovymi koeficientami. Iz togo kako parametrično predstavjenje dostavajemo prosto jednačenje s neopreděljenymi koeficientami. Kako od jego prějdti k libovoljno mnogim jednačenjam ``bez afekta'' nad vsakym čislovym poljem, pokazuje Hilbertova teorema o nerazložimosti.

Nyně iz teoremov Lürotha i Castelnuova neposrědnje slěduje eksistencija minimalnoj bazy za vse grupy, ktore nastupajut pri jednačenjah 3-go i 4-go stepena; pri tom se dalje pokazuje, že ta minimalna baza imaje racionalne čislove koeficienty. Tako možno za vsako libovoljno čislovo polje racionalno konstruovati cělost jednačenj 3-go i 4-go stepena s prědpisanoju grupoju. Za ciklične i diedralne grupy eksistuje minimalna baza s poljem $n$-yh korenjev iz jedinice kako koeficientnym područjem; to samo važi za něktore dalje metaciklične grupy. Pytanje, či minimalna baza vobće jest harakterna za něktore grupy, mora ostati nerěšeno.

3. Da dojdemo k poslědnjemu bazovomu pojmu, razsmatrjam cělost polinomov soderžanyh v polju. Pod ``integralnoju bazoju'' razuměju konečno čislo tyh polinomov tako, že vsak polinom polja može se predstaviti kako cěla racionalna kombinacija togo konečnogo čisla polinomov, s koeficientami iz danogo koeficientnogo područja. Pytanje o eksistenciji integralnoj bazy -- ktore, napr., kako specialny slučaj vključa konečnost sistema invariantov -- bylo postavjeno Hilbertom v jego ``Matematičnyh problemah''\footnote{D. Hilbert: Mathematische Probleme. Doklad, Paris 1900. Problem 14. Göttinger Nachr. 1900.} v malo drugoj formulaciji, kako problem relativno cělyh funkcij.

Za nyně mogu ukazati samo jednu klasu polj, za ktore eksistuje integralna baza. Imenno, ako polje soderži sistem iz $n$ polinomov tako, že homogena rezultanta členov najvyššego stepena tyh polinomov jest različna od nula, togda eksistuje integralna baza; ona jest dana koeficientami vseh $u$ v jednačenju, nerazložimom nad poljem, za linearnu formu:
\[
  u_0=u_1x_1+u_2x_2+\cdots+u_nx_n.\footnotemark
\]
\footnotetext{To nerazložimo jednačenje v slučaju jednoj neopreděljenej nahodi se u E. Steinitza v dokazu Lürothovej teoremy.}
Ako osoblivo vrědnost rezultante jest ravna jedinici koeficientnogo područja, togda jest osigurana takože integralnost predstavjenja. Priměr toj klasy polj davajut Lagrangeve rodove oblasti, bo rezultanta elementarnyh simetričnyh funkcij imaje vrědnost $\pm 1$. Daljny priměr (bez integralnosti) jest dany vsimi poljami, ktore soderžat samo jedin algebraično nezavisny polinom; integralna baza tu jest identična s Lürothovoju funkcijeju najmanjšego podpolja, ktore soderži vse polinomy. Tako, kako jest znajemo, vse cěle racionalne projektivne invarianty kvadratičnoj formy od $n$ prěmennyh sut cěle funkcije diskriminanta.

Dano uslovje jest samo dostatočno, nikako ne nužno za eksistenciju integralnoj bazy. Možno legko konstruovati polja, pri ktoryh rezultanta vsakih $n$ polinomov izčezaje, i ktore vseže imajut integralnu bazu danu koeficientami togo nerazložimogo jednačenja. Nastavaje pytanje, či možda i v najobćejšem slučaju koeficienty togo jednačenja davajut integralnu bazu.

\end{document}
